%%═══════════════════════════════════════════════════════════════════
%% Lecture 08 — Fermion Bilinears, Renormalizability, and the Road
%%              to the Standard Model
%% 77609 — Introduction to Elementary Particles
%% Date: 03/05/2026
%% Lecturer: Prof. Yonit Hochberg
%%═══════════════════════════════════════════════════════════════════

\documentclass[../main.tex]{subfiles}
\usepackage{preamble}

\begin{document}

\section{Lecture 08 — Fermion Bilinears, Renormalizability, and the Road to the Standard Model}\label{sec:lec08}
\begin{center}
\begin{tabular}{ll}
  \textbf{Date:}\quad 03/05/2026 & \\
\end{tabular}
\end{center}
\hrule\vspace{1.5em}
%%──────────────────────────────────────────────────────────────────
\subsection*{Overview}

This lecture completes the foundational treatment of Dirac fermions and introduces the guiding principles behind the Standard Model Lagrangian.

\begin{enumerate}
  \item \textbf{Projection operators}: extracting $\psi_L$ and $\psi_R$ from the Dirac spinor via $P_{L,R}$.
  \item \textbf{Dirac equation in Weyl form}: coupling of chiralities by the mass; massless decoupling into Weyl equations.
  \item \textbf{Fermion bilinears}: the complete set of Lorentz-covariant objects bilinear in $\psi$; scalar, vector, pseudo-vector, pseudo-scalar, and tensor.
  \item \textbf{Feynman slash notation}: $\slashed{A} = \gamma^\mu A_\mu$; compact form of the Dirac equation and Lagrangian.
  \item \textbf{Component notation}: spinor index structure of the Dirac Lagrangian; distinguishing spinor indices from Lorentz indices.
  \item \textbf{Mass dimensions}: $[\psi] = \tfrac{3}{2}$, $[\phi] = 1$; renormalizability limits interactions to at most two fermion fields.
  \item \textbf{Dirac mass term}: structure $m\bar\psi\psi = m(\bar\psi_L\psi_R + \bar\psi_R\psi_L)$; Majorana mass briefly mentioned.
  \item \textbf{Fermion propagator}: $i(\slashed{p}+m)/(p^2-m^2)$; comparison with the scalar propagator.
  \item \textbf{Chirality and parity}: handedness is reversed under parity.
  \item \textbf{Requirements on the SM Lagrangian}: locality, translational invariance, Poincaré symmetry, analyticity, renormalizability, naturalness, internal symmetries.
  \item \textbf{The Standard Model}: the statement $SU(3)\times SU(2)\times U(1)$ spontaneously broken gauge symmetry accounts for almost all experimental data.
\end{enumerate}

%%──────────────────────────────────────────────────────────────────
\subsection*{Projection Operators and Chirality Decomposition}

\subsubsection*{Physical Motivation}

The Dirac spinor $\psi$ is a reducible representation of the Lorentz group: it decomposes into a left-handed and a right-handed two-component Weyl spinor, each of which transforms irreducibly.  To isolate each piece we need a projector that commutes with Lorentz transformations.  The matrix $\gamma^5$ (whose eigenvalues $\pm1$ distinguish the two chiralities) provides exactly such a projector.

\subsubsection*{Definition}

Given the $\gamma^5$ matrix (defined as $\gamma^5 \equiv i\gamma^0\gamma^1\gamma^2\gamma^3$, with $(\gamma^5)^2 = \mathbf{1}$), the projection operators onto left- and right-handed Weyl subspaces are:
\begin{equation}
  \boxed{P_L \equiv \frac{1-\gamma^5}{2}, \qquad P_R \equiv \frac{1+\gamma^5}{2}.}
  \label{eq:lec08:projectors}
\end{equation}
In the standard (Weyl/chiral) basis,
\[
  \gamma^5 = \begin{pmatrix}-\mathbf{1}&0\\0&\mathbf{1}\end{pmatrix},
\]
so
\[
  P_L = \begin{pmatrix}\mathbf{1}&0\\0&0\end{pmatrix}, \qquad
  P_R = \begin{pmatrix}0&0\\0&\mathbf{1}\end{pmatrix}.
\]

\textit{Physical significance:} $P_L$ and $P_R$ extract, respectively, the top two and bottom two components of $\psi$ — precisely the left-handed and right-handed Weyl spinors $\psi_L$ and $\psi_R$ introduced in Lecture~07.

\subsubsection*{Properties}

The projectors satisfy the standard idempotency and orthogonality relations:
\begin{equation}
  P_L^2 = P_L, \quad P_R^2 = P_R, \quad P_L P_R = P_R P_L = 0, \quad P_L + P_R = \mathbf{1}.
  \label{eq:lec08:proj-props}
\end{equation}
Acting on the four-component Dirac spinor:
\begin{equation}
  \psi_L \equiv P_L\psi = \begin{pmatrix}\psi_L\\0\end{pmatrix}, \qquad
  \psi_R \equiv P_R\psi = \begin{pmatrix}0\\\psi_R\end{pmatrix}.
  \label{eq:lec08:weyl-proj}
\end{equation}

%%──────────────────────────────────────────────────────────────────
\subsection*{Dirac Equation in Weyl Component Form}

\subsubsection*{Physical Motivation}

Writing the Dirac equation in terms of $\psi_L$ and $\psi_R$ reveals how the mass term couples the two chiralities.  In the massless limit the equation decouples into two independent equations — the Weyl equations — one for each chirality.  This decomposition is physically important because the weak force acts \emph{only on $\psi_L$}, so the Weyl basis is the natural language for the Standard Model.

\subsubsection*{Matrix form}

In the Weyl basis, writing $\psi = (\psi_L, \psi_R)^T$, the Dirac equation $(i\gamma^\mu\partial_\mu - m)\psi = 0$ becomes:
\begin{equation}
  \begin{pmatrix}
    -m & i\sigma^\mu\partial_\mu \\
    i\bar\sigma^\mu\partial_\mu & -m
  \end{pmatrix}
  \begin{pmatrix}\psi_L \\ \psi_R\end{pmatrix} = 0,
  \label{eq:lec08:dirac-weyl-matrix}
\end{equation}
where $\sigma^\mu = (\mathbf{1},\,\sigma^i)$ and $\bar\sigma^\mu = (\mathbf{1},\,-\sigma^i)$ (recall Lecture~07).  Explicitly, this gives two coupled equations:
\begin{align}
  i\bar\sigma^\mu\partial_\mu\,\psi_L &= m\,\psi_R, \label{eq:lec08:dirac-L}\\
  i\sigma^\mu\partial_\mu\,\psi_R     &= m\,\psi_L. \label{eq:lec08:dirac-R}
\end{align}
The mass $m$ couples $\psi_L$ and $\psi_R$: the kinetic part drives each chirality forward in spacetime while the mass ``flips'' the chirality.

\subsubsection*{Massless limit: Weyl equations}

When $m = 0$, equations \eqref{eq:lec08:dirac-L} and \eqref{eq:lec08:dirac-R} decouple completely:
\begin{equation}
  \boxed{
  i\bar\sigma^\mu\partial_\mu\,\psi_L = 0, \qquad i\sigma^\mu\partial_\mu\,\psi_R = 0.
  }\label{eq:lec08:weyl-equations}
\end{equation}
Explicitly
\begin{equation}
    \begin{cases}
i \left[ I \cdot \frac{\partial}{\partial t} + \sum_{i=1}^3 \sigma^i \frac{\partial}{\partial x^i} \right] \psi_R = 0 & \text{(Massless Right-Handed Weyl Equation)} \\
i \left[ I \cdot \frac{\partial}{\partial t} - \sum_{i=1}^3 \sigma^i \frac{\partial}{\partial x^i} \right] \psi_L = 0 & \text{(Massless Left-Handed Weyl Equation)}
\end{cases}
\end{equation}
These are the \textbf{Weyl equations}.  Each is an equation for a single two-component (Weyl) spinor.

\textit{Physical significance:} Neutrinos in the original Standard Model were modeled as \emph{massless} Weyl fermions; the discovery of neutrino masses is therefore a sign of physics beyond the SM.  The Weyl equations are the minimal relativistic wave equations for massless spin-$\tfrac{1}{2}$ particles.

%%──────────────────────────────────────────────────────────────────
\subsection*{Fermion Bilinears and Their Lorentz Properties}

\subsubsection*{Physical Motivation}

When constructing interaction terms in a Lagrangian, we need Lorentz-invariant combinations of fermion fields.  A single $\psi$ transforms non-trivially; only specific bilinears $\bar\psi\,\Gamma\,\psi$ (where $\Gamma$ is a combination of $\gamma$ matrices) have definite Lorentz transformation properties.  Cataloguing all such bilinears gives the complete menu of fermion interactions allowed by Lorentz symmetry.

\subsubsection*{The complete basis}

The set $\{\mathbf{1},\,\gamma^\mu,\,\gamma^5,\,\gamma^\mu\gamma^5,\,\Sigma^{\mu\nu}\}$ spans the full space of $4\times4$ matrices (16 independent elements), where the \textbf{rank-2 antisymmetric} tensor is:
\begin{equation}
  \Sigma^{\mu\nu} \equiv \frac{i}{2}\bigl[\gamma^\mu,\,\gamma^\nu\bigr]
  = \frac{i}{2}(\gamma^\mu\gamma^\nu - \gamma^\nu\gamma^\mu).
  \label{eq:lec08:Sigma-def}
\end{equation}
Note $\Sigma^{\mu\nu}$ is antisymmetric: $\Sigma^{\mu\nu} = -\Sigma^{\nu\mu}$, so it has $\binom{4}{2} = 6$ independent components.

The bilinears and their Lorentz transformation properties are summarized in the following table:

\begin{center}
\renewcommand{\arraystretch}{1.4}
\begin{tabular}{llcc}
  \toprule
  Bilinear & Lorentz type & Components & Notes \\
  \midrule
  $\bar\psi\psi$ & Scalar & 1 & Invariant under all Lorentz transf. \\
  $\bar\psi\gamma^\mu\psi$ & Vector & 4 & Transforms as $V^\mu \to \Lambda^\mu_{\ \nu}V^\nu$ \\
  $\bar\psi\gamma^5\psi$ & Pseudo-scalar & 1 & Extra $-$ sign under parity \\
  $\bar\psi\gamma^\mu\gamma^5\psi$ & Pseudo-vector (axial) & 4 & Extra $-$ sign under parity \\
  $\bar\psi\Sigma^{\mu\nu}\psi$ & Tensor & 6 & Antisymmetric in $\mu\nu$ \\
  \bottomrule
\end{tabular}
\end{center}

Total: $1+4+1+4+6 = 16$ independent components, forming a complete basis for $4\times4$ matrices. \checkmark

\begin{remark}[Pseudo- means extra minus sign under parity]
The prefix ``pseudo'' indicates that the object transforms like its non-pseudo counterpart under proper Lorentz transformations (rotations and boosts), but acquires an \emph{additional minus sign} under the parity transformation $\bvec{x}\to -\bvec{x}$.  For example, $\bar\psi\gamma^\mu\gamma^5\psi$ is called an \textbf{axial vector}: it transforms as a vector under boosts and rotations, but changes sign under parity, unlike the vector $\bar\psi\gamma^\mu\psi$.
\end{remark}

\textit{Physical significance:} The vector bilinear $\bar\psi\gamma^\mu\psi$ is the electromagnetic current $j^\mu_{\rm EM}$.  The axial-vector bilinear $\bar\psi\gamma^\mu\gamma^5\psi$ appears in weak interactions; the combination $V - A$ (vector minus axial) is the hallmark of parity violation in the weak force, which is one of the most surprising experimental facts about Nature.

%%──────────────────────────────────────────────────────────────────
\subsection*{Feynman Slash Notation}

A compact notation for contracting a four-vector with the gamma matrices:
\begin{equation}
  \boxed{{\displaystyle {A\!\!\!/}\ {\stackrel {\mathrm {def} }{=}}\ \gamma ^{0}A_{0}+\gamma ^{1}A_{1}+\gamma ^{2}A_{2}+\gamma ^{3}A_{3}}}
  \label{eq:lec08:slash-def}
\end{equation}
This applies to any object with a Lorentz index $\mu$.  For example:
\begin{itemize}
  \item $\slashed{\partial} = \gamma^\mu\partial_\mu$ (derivative operator).
  \item $\slashed{p} = \gamma^\mu p_\mu$ (four-momentum on a fermion line).
  \item $\slashed{A} = \gamma^\mu A_\mu$ (gauge boson field $A^\mu$, relevant for QED).
\end{itemize}

The Dirac equation and the Dirac Lagrangian in slash notation:
\begin{equation}
  (i\slashed\partial - m)\psi = 0,
  \label{eq:lec08:dirac-slash}
\end{equation}
\begin{equation}
  \boxed{\lagr_{\rm Dirac} = \bar\psi\,(i\slashed\partial - m)\,\psi.}
  \label{eq:lec08:Lagr-slash}
\end{equation}
The kinetic term is $\bar\psi\,i\slashed\partial\,\psi$ and the mass term is $-m\bar\psi\psi$.  Both are Lorentz invariant: $\bar\psi\psi$ is a Lorentz scalar (from the bilinear table above), and $\bar\psi\gamma^\mu\partial_\mu\psi$ is the contraction of a Lorentz vector with a four-derivative, which is also a scalar.

%%──────────────────────────────────────────────────────────────────
\subsection*{Component Notation for the Dirac Lagrangian}

\subsubsection*{Physical Motivation}

It is useful to write the Dirac Lagrangian with explicit spinor (component) indices to understand what kind of object $\gamma^\mu$ is and what indices it carries.  A common source of confusion is conflating the spinor index of $\psi$ (which runs 1 to 4 over spin and particle/antiparticle degrees of freedom) with the Lorentz index $\mu$ (which runs 0 to 3 over spacetime directions).

\subsubsection*{Explicit index structure}

Let $\alpha, \beta = 1,2,3,4$ label spinor components.  Then:
\begin{equation}
  \lagr_{\rm Dirac}
  = \bar\psi_\alpha\Bigl[i\,(\gamma^\mu)_{\alpha\beta}\,\partial_\mu
    - m\,\delta_{\alpha\beta}\Bigr]\psi_\beta,
  \label{eq:lec08:Lagr-components}
\end{equation}
where $(\gamma^\mu)_{\alpha\beta}$ is the $(\alpha,\beta)$ matrix element of the $4\times4$ matrix $\gamma^\mu$.  Key points:
\begin{itemize}
  \item $\psi_\beta$ is a four-component column vector (spinor), with index $\beta$ running over spinor space.
  \item $(\gamma^\mu)_{\alpha\beta}$ is a $4\times4$ matrix acting on spinor space, and $\mu$ is its Lorentz index — not a spinor index.
  \item $\partial_\mu$ carries only a Lorentz index; it acts on the spacetime argument of $\psi$, independent of the spinor index.
  \item The spinor indices $\alpha$ and $\beta$ are summed over (Einstein convention on spinor indices), just like the Lorentz index $\mu$ is summed.
\end{itemize}

\begin{remark}[Two separate ``fours'']
The Dirac spinor has four components, but these four have nothing to do with the four spacetime directions.  One set of ``four'' lives in spinor space; the other lives in Minkowski space.  Keeping the two distinct prevents errors in Lorentz transformation arguments: a Lorentz boost does not mix the spinor-$\alpha$ index, it mixes the $\mu$ index.
\end{remark}

%%──────────────────────────────────────────────────────────────────
\subsection*{Mass Dimensions of Fermionic Fields}

\subsubsection*{Physical Motivation}

In natural units the action $S = \int d^4x\,\lagr$ must be dimensionless (it enters $e^{iS}$).  Since $[d^4x] = [{\rm mass}]^{-4}$, the Lagrangian density must satisfy $[\lagr] = [{\rm mass}]^4$.  Reading off the dimension of $\psi$ from the Lagrangian then constrains what interaction terms are allowed: this is the origin of \emph{renormalizability}.

\subsubsection*{Derivation of $[\psi] = 3/2$}

From the kinetic term $\bar\psi\,i\slashed\partial\,\psi$:
\[
  [\bar\psi\,\slashed\partial\,\psi]
  = [\psi]^2\,[\partial]
  = [\psi]^2 \cdot 1
  \stackrel{!}{=} 4,
\]
so $[\psi]^2 = 3$ and:
\begin{equation}
  \boxed{[\psi] = \tfrac{3}{2}.}
  \label{eq:lec08:psi-dim}
\end{equation}
Independently, from the mass term $m\bar\psi\psi$: $[m][\psi]^2 = 1\cdot 3 = 4$. \checkmark

For comparison, the scalar kinetic term $(\partial\phi)^2$ gives $[\phi]^2 \cdot 1^2 = 4$, so:
\begin{equation}
  [\phi] = 1.
  \label{eq:lec08:phi-dim}
\end{equation}

\textit{Physical significance:} Fermionic fields have a higher mass dimension ($3/2$) than scalar fields ($1$).  This makes fermion interactions more constrained: each fermion ``uses up'' more of the dimension budget.

%%──────────────────────────────────────────────────────────────────
\subsection*{Renormalizability and the Two-Fermion Constraint}

\subsubsection*{Physical Motivation}

A fundamental principle is that the full theory of nature should be \textbf{renormalizable}: all interaction terms in the Lagrangian should have mass dimension $\leq 4$.  (Terms of dimension $> 4$ come with inverse powers of some large mass scale $\Lambda$, signaling that the theory is only an effective description valid at energies $\ll \Lambda$.)  With $[\psi] = 3/2$, this severely limits how many fermion fields can appear in one interaction term.

\subsubsection*{The constraint}

A renormalizable Lagrangian contains only terms with mass dimension at most 4:
\begin{equation}
  \boxed{\text{Renormalizable: all terms have } [\text{fields}+\text{derivatives}] \leq 4.}
  \label{eq:lec08:renorm}
\end{equation}

For fermions, two fields already contribute $2 \times \tfrac{3}{2} = 3$ to the dimension.  A third fermion would require dimension $3 + \tfrac{3}{2} = \tfrac{9}{2} > 4$.  Therefore:

\begin{equation}
  \boxed{\text{At most two fermion fields per term in a renormalizable Lagrangian.}}
  \label{eq:lec08:two-fermion}
\end{equation}

Moreover, a single fermion cannot appear alone (it would be odd under Lorentz transformation and not Lorentz invariant), so fermions always enter in pairs: either $\bar\psi\psi$ (dim 3) or $\bar\psi\psi$ multiplied by a scalar field or a single derivative (up to dim 4).  Possible renormalizable terms involving two fermions include:
\begin{itemize}
  \item $m\bar\psi\psi$ — mass term (dim 4 with $m$ of dim 1).
  \item $y\,\phi\,\bar\psi\psi$ — Yukawa coupling to scalar $\phi$ (dim $1+3=4$, coupling $y$ dimensionless).
  \item $\bar\psi\,i\slashed\partial\,\psi$ — kinetic term (dim $3/2 + 1 + 3/2 = 4$).
\end{itemize}

\textit{Physical significance:} This constraint is why the Standard Model Lagrangian has its specific form.  Four-fermion interactions (Fermi theory of weak decays) have dim 6, signaling that Fermi theory is an effective theory — valid at low energies but breaking down near the $W$ boson mass, where the full renormalizable SM takes over.

%%──────────────────────────────────────────────────────────────────
\subsection*{Structure of the Dirac Mass Term}

\subsubsection*{Physical Motivation}

We saw from the Weyl-component form of the Dirac equation that the mass term couples $\psi_L$ and $\psi_R$.  Let us verify this explicitly and understand why this is inevitable for a Lorentz-invariant mass term.

\subsubsection*{Chirality structure of $m\bar\psi\psi$}

Using $\psi = \psi_L + \psi_R$ (with $\psi_L = P_L\psi$, $\psi_R = P_R\psi$) and $\bar\psi = \psi^\dagger\gamma^0$:
\begin{align}
  \bar\psi\psi
  &= (\bar\psi_L + \bar\psi_R)(\psi_L + \psi_R) \notag\\
  &= \bar\psi_L\psi_L + \bar\psi_R\psi_R + \bar\psi_L\psi_R + \bar\psi_R\psi_L.
\end{align}
One can show (using the $\gamma^5$ eigenvalue properties) that $\bar\psi_L\psi_L = \bar\psi_R\psi_R = 0$, so:
\begin{equation}
  \boxed{m\bar\psi\psi = m\bigl(\bar\psi_L\psi_R + \bar\psi_R\psi_L\bigr).}
  \label{eq:lec08:dirac-mass}
\end{equation}
This is the \textbf{Dirac mass term}: it necessarily couples left- and right-handed fields.

\textit{Physical significance:} A Dirac mass term requires both a left-handed and a right-handed component to exist.  In the SM, the left-handed quarks and leptons carry different electroweak quantum numbers than their right-handed counterparts.  Mass terms that couple $L$ and $R$ therefore break the electroweak symmetry — which is why fermion masses in the SM can only arise after electroweak symmetry breaking (the Higgs mechanism).

\begin{remark}[Majorana mass]
A \textbf{Majorana mass term} involves only one chirality (e.g., only $\psi_L$ or only $\psi_R$).  Such a term is allowed for electrically neutral particles (e.g., neutrinos) because the charge conjugate of $\psi_L$ can serve as the missing chirality.  Majorana masses are not covered in detail in this course.
\end{remark}

%%──────────────────────────────────────────────────────────────────
\subsection*{Fermion Propagator}

\subsubsection*{Physical Motivation}

In Feynman diagram computations, every internal fermion line contributes a propagator.  The fermion propagator is obtained from the free Dirac Lagrangian in the same way as the scalar propagator, but its numerator carries the momentum $\slashed{p}+m$ because of the first-order nature of the Dirac equation.

\subsubsection*{Result}

The Feynman propagator for a Dirac fermion with momentum $p$ and mass $m$ is:
\begin{equation}
  \boxed{\langle\psi\bar\psi\rangle_F = \frac{i(\slashed{p}+m)}{p^2 - m^2}.}
  \label{eq:lec08:fermion-prop}
\end{equation}
In a Wick contraction:
\[
  \overbrace{\psi\cdots\bar\psi}^{\phantom{X}} = \frac{i(\slashed{p}+m)}{p^2-m^2}.
\]
Compare with the scalar propagator:
\[
  \text{Scalar:}\quad \frac{i}{p^2-m^2}; \qquad
  \text{Fermion:}\quad \frac{i(\slashed{p}+m)}{p^2-m^2}.
\]

\textit{Physical significance:} The extra factor $(\slashed{p}+m)$ in the numerator means the fermion propagator falls off as $1/p$ for large $p$, compared to $1/p^2$ for the scalar.  This different UV behavior leads to different divergence structures for loop diagrams involving fermions vs.\ scalars, with important consequences for renormalization.

%%──────────────────────────────────────────────────────────────────
\subsection*{Chirality and Parity}

The \textbf{chirality} (or \textbf{handedness}) of a Weyl spinor is its $\gamma^5$ eigenvalue: $+1$ for $\psi_R$, $-1$ for $\psi_L$.

Under a \textbf{parity transformation} ($\bvec{x} \to -\bvec{x}$, $t\to t$):
\begin{equation}
  \boxed{\psi_L \;\xrightarrow{\text{parity}}\; \psi_R, \qquad \psi_R \;\xrightarrow{\text{parity}}\; \psi_L.}
  \label{eq:lec08:chirality-parity}
\end{equation}
Chirality is reversed under parity.  This is why a theory that treats $\psi_L$ and $\psi_R$ differently (as the SM does for the weak interaction) \emph{violates parity} — an experimentally confirmed fact (Wu experiment, 1956).

%%──────────────────────────────────────────────────────────────────
\subsection*{Requirements on the Lagrangian: Toward the Standard Model}

\subsubsection*{Physical Motivation}

We are now ready to ask: what is the Lagrangian that describes nature?  Rather than guessing, we impose a set of physical requirements that any Lagrangian must satisfy.  Each requirement restricts the allowed terms and guides us toward the Standard Model.

\subsubsection*{The requirements}

The Lagrangian $\lagr[\phi_i, \partial_\mu\phi_i]$ must satisfy:

\begin{enumerate}
  \item \textbf{Locality and analyticity}: $\lagr$ depends on fields and their derivatives \emph{at a single spacetime point} $x^\mu$.  This defines a \textbf{local field theory}.  Moreover, $\lagr$ should be analytic (polynomial) in the fields — a necessary condition for perturbation theory to apply.

  \item \textbf{Translational invariance}: $\lagr$ does not depend \emph{explicitly} on $x^\mu$ (only implicitly through the fields).  By Noether's theorem, this guarantees conservation of energy and momentum.

  \item \textbf{Poincaré invariance}: $\lagr$ must be invariant under the full \textbf{Poincaré group} = (Lorentz transformations) + (spacetime translations).  This is the fundamental spacetime symmetry of special relativity.  It is Poincaré symmetry — not just Lorentz symmetry — that underlies the concepts of \textbf{mass} and \textbf{spin}: these quantum numbers label the irreducible representations of the Poincaré group.

  \item \textbf{Renormalizability}: $\lagr$ contains only terms with mass dimension $\leq 4$.  Higher-dimensional terms signal an effective theory with a UV cutoff $\Lambda$; a truly fundamental theory should be renormalizable.  (Non-renormalizable terms are suppressed by powers of $E/\Lambda \ll 1$.)

  \item \textbf{Naturalness}: Every term in $\lagr$ that is \emph{not forbidden} by a symmetry \emph{should} appear.  Stated differently, if symmetries allow a term, it must be present — one should not arbitrarily set coefficients to zero.

  \item \textbf{Internal symmetry groups}: Beyond spacetime symmetry, one may require $\lagr$ to be invariant under global or local internal symmetry groups (e.g., $U(1)$, $SU(2)$, $SU(3)$).  By Noether's theorem, each symmetry implies a conserved current and conserved charge.

\end{enumerate}

\subsubsection*{Summary of constraints}

\begin{center}
\renewcommand{\arraystretch}{1.4}
\begin{tabular}{ll}
  \toprule
  Requirement & Consequence \\
  \midrule
  Locality & Fields interact at a point; no action at a distance \\
  Translational invariance & Energy-momentum conservation \\
  Lorentz (Poincaré) invariance & Well-defined mass and spin \\
  Renormalizability & Dim $\leq4$ terms only; at most 2 fermions per term \\
  Naturalness & No fine-tuned cancellations without a symmetry reason \\
  Internal symmetries & Conserved charges; gauge interactions \\
  \bottomrule
\end{tabular}
\end{center}

%%──────────────────────────────────────────────────────────────────
\subsection*{The Standard Model}

Applying all of the requirements above, one arrives at a remarkably constrained set of possible Lagrangians.  The verdict from over 50 years of experiment is:

\begin{center}
\fbox{\parbox{0.85\linewidth}{
\textbf{Statement of the Standard Model:}\\[4pt]
Almost all experimental data for elementary particles and their interactions can be explained by a \textbf{spontaneously broken} $SU(3)\times SU(2)\times U(1)$ \textbf{gauge theory}.
}}
\end{center}

Each factor of the gauge group has a physical meaning:
\begin{itemize}
  \item $SU(3)$ — the strong force ($\QCD$); gluons are the gauge bosons; quarks carry colour charge.
  \item $SU(2)\times U(1)$ — the electroweak force; the $W^\pm$, $Z^0$, and photon are the gauge bosons; spontaneous symmetry breaking by the Higgs field gives them mass.
\end{itemize}

The word \emph{almost} is important: the Standard Model does not account for dark matter, dark energy, gravity, neutrino masses (in its minimal form), or the matter-antimatter asymmetry.  The outstanding goal of modern particle physics is to find experimental deviations from $\SM{}$ predictions and identify the more fundamental theory.

\textit{Physical significance:} The rest of the course will be devoted to understanding each word in this statement — what ``spontaneous symmetry breaking'' means, how gauge symmetries generate interactions, and why this particular group $SU(3)\times SU(2)\times U(1)$ describes nature.

%%──────────────────────────────────────────────────────────────────
\subsection*{To Remember}

\begin{itemize}
  \item \textbf{Projection operators}: $P_L = \tfrac{1}{2}(1-\gamma^5)$, $P_R = \tfrac{1}{2}(1+\gamma^5)$; $\psi_{L,R} = P_{L,R}\psi$.

  \item \textbf{Weyl equations} (massless limit): $i\bar\sigma^\mu\partial_\mu\psi_L = 0$, $i\sigma^\mu\partial_\mu\psi_R = 0$.

  \item \textbf{Dirac mass couples chiralities}: $m\bar\psi\psi = m(\bar\psi_L\psi_R + \bar\psi_R\psi_L)$.

  \item \textbf{Fermion bilinears}: scalar ($\bar\psi\psi$), vector ($\bar\psi\gamma^\mu\psi$), pseudo-scalar ($\bar\psi\gamma^5\psi$), axial vector ($\bar\psi\gamma^\mu\gamma^5\psi$), tensor ($\bar\psi\Sigma^{\mu\nu}\psi$ with 6 components).

  \item \textbf{Pseudo} = extra $-$ sign under parity; \textbf{axial} = pseudo-vector.

  \item \textbf{Slash notation}: $\slashed{A} = \gamma^\mu A_\mu$; Dirac eq.\ $(i\slashed\partial - m)\psi = 0$.

  \item \textbf{Mass dimensions}: $[\psi] = \tfrac{3}{2}$, $[\phi] = 1$, $[\lagr] = 4$.

  \item \textbf{Renormalizability} $\Rightarrow$ at most two fermion fields per term.

  \item \textbf{Fermion propagator}: $i(\slashed{p}+m)/(p^2-m^2)$; falls as $1/p$, not $1/p^2$.

  \item \textbf{Chirality reverses under parity}: $\psi_L \leftrightarrow \psi_R$.

  \item \textbf{Requirements on $\lagr$}: locality, translational invariance, Poincaré invariance, analyticity, renormalizability, naturalness, internal symmetries.

  \item \textbf{The Standard Model}: spontaneously broken $SU(3)\times SU(2)\times U(1)$ gauge theory accounts for almost all particle physics data.
\end{itemize}

%%──────────────────────────────────────────────────────────────────
\subsection*{Recommended Reading}

\begin{itemize}
  \item \textbf{Peskin \& Schroeder, Chapter~3} — Projection operators, chirality, Weyl spinors, fermion bilinears, and the Dirac propagator.  The reference explicitly used in this course for fermion material.

  \item \textbf{Tong, \textit{Lectures on Particle Physics}, Chapter~1 (draft)}
    (\href{https://www.damtp.cam.ac.uk/user/tong/pp.html}{\texttt{damtp.cam.ac.uk/user/tong/pp.html}}):
    Chapter~1 covers the requirements on the SM Lagrangian (locality, renormalizability, gauge symmetry) and states the SM gauge group — directly parallel to the second half of this lecture.

  \item \textbf{Halzen \& Martin, Chapter~5} — Fermion bilinears, parity, the $V-A$ structure of weak interactions; accessible treatment.

  \item \textbf{Griffiths, Chapter~9} — Introduction to the Standard Model gauge group; complements the lecture's closing statement.
\end{itemize}

\end{document}
